\section{3D Printing}

Manufacturing processes are commonly divided into two main categories: additive manufacturing and subtractive manufacturing. Subtractive manufacturing creates objects by removing material from a larger block using methods such as cutting, drilling, or milling. Traditional machining techniques like CNC milling are typical examples of this approach. In contrast, additive manufacturing builds objects by gradually adding material layer by layer according to a digital model. This method allows manufacturers to produce complex shapes and structures that would be difficult or impossible to achieve using traditional subtractive techniques.

3D printing is one of the most well-known forms of additive manufacturing. In this process, a digital 3D model is first designed using computer-aided design (CAD) software. The 3D printer then creates the object by depositing or solidifying material layer by layer until the entire shape is formed. Materials used in 3D printing can include plastics, resins, metals, and even biological materials. Because of its layer-by-layer construction, 3D printing enables rapid prototyping and allows designers to easily modify and test new ideas.

3D printing is particularly strong in rapid prototyping, customization, and producing complex geometries with minimal material waste. It allows designers and engineers to quickly turn digital designs into physical models, making the development process faster and more flexible. However, 3D printing is not always suitable for large-scale mass production because it is generally slower than traditional manufacturing methods. In addition, some 3D-printed materials may not provide the same strength, durability, or surface finish as parts produced using conventional manufacturing techniques.